\section{The original fraction-field norm of an arbitrary DVR}

\begin{theorem}[The actual DVR norm extends to its actual fraction field]
\label{thm:actual-dvr-fraction-point}
\lean{EnvelopeProperty.Berk.dvrFractionPoint, EnvelopeProperty.Berk.dvrFractionPoint_val_algebraMap, EnvelopeProperty.Berk.dvrFractionPoint_val_algebraMap_le_one, EnvelopeProperty.Berk.dvrFractionPoint_integers}
\leanok
Let $V$ be a DVR over the trivially valued field $k$, let $K$ be its actual
fraction field, and let $c>0$. The original norm
$|a|=\exp(-c\operatorname{ord}_V a)$ extends to the actual field $K$,
with its original values on every element of $V$. Its ring of integers is
exactly the original $V$.
\end{theorem}
\begin{proof}
The original norm is nonzero on every nonzero element of $V$. Apply the
actual localization extension to the original fraction field; its value on
$a/b$ is $|a|/|b|$. Field coefficients remain trivial by the original scalar
tower. Every original element of $V$ has norm at most one. Conversely suppose
$x$ has norm at most one but is not in $V$. The valuation-ring property of
the original DVR gives $x^{-1}=a\in V$. The element $a$ is a nonunit, since
otherwise $x$ would also lie in $V$. Therefore its order is positive and
$|a|<1$, whereas $|x|\le1$ and $xa=1$ imply $|a|\ge1$, a contradiction.
\end{proof}
